\section{Ten positive unknowns with an 89-operation evaluator}
\label{sec:positive}

The direct result above keeps modest coefficients and unrestricted integer
unknowns. A complementary calculation shortens the existing compact native
route. This section's domain is explicitly $\Z_{>0}^{10}$, and its large
fixed coefficients are inherited from the coherent compact44 BFS tuple.
It is not a 10-variable unrestricted-integer result.

\begin{center}
\begin{tabular}{lrrr}
\toprule
Polynomial & Variables & Operations & Exact degree\\
\midrule
Uncentered & 10 & 89 & 288\\
Upper center & 10 & 91 & 283\\
Sharp center & 10 & 92 & 278\\
Direct, uncentered & 10 & 94 & 262\\
Direct, upper center & 10 & 96 & 258\\
Direct, sharp center & 10 & 97 & 254\\
\bottomrule
\end{tabular}
\end{center}

The centered parents in the incoming saturated-Pell packet cost, in the
same order, $92/283$, $93/278$, $98/258$, and $99/254$, where each pair is
operations/degree \cite{saturated}. Thus all four centered companions
improve cost at unchanged degree. The uncentered 89-operation equation is
a new cost/degree tradeoff. The earlier square-gap and auxiliary-index
ideas in the incoming structural-compression packet are explicitly credited
\cite{structural}; the new step is their combination with coprime Pell
fusion and the resulting replacement of the transport majorant.

\subsection{Complete polynomial definition}
Use the positive unknowns
\[
 J,F,\alpha,i,T,w,Y,n,\eta,\zeta.
\]
On the direct chart replace $w$ by a positive unknown $x$. The four
constants $\beta,K_0,M_C,M_F$ are the entire compact44 BFS tuple supplied
in \file{positive/coefficients/}. In particular $\beta=2^{590069}-1$.
The full fixed finite recipe, including its 436-character seed and all
finite index bounds, is \file{positive/coefficients/recipe_compact44.py}.
It reconstructs every coefficient and checks its complete raw-byte hash.
The high mask $M_F$ is already shifted; no second shift or mixture with
the differently ordered original44 tuple is permitted.

Define only polynomial abbreviations:
\begin{align*}
 q&=\beta J+1,& \mu&=q^2+1,\\
 k&=\eta+\zeta,&c&=kY+\eta,\\
 X&=wq\quad\text{(ordinary)},&&
 X=q+x\quad\text{(direct)},\\
 A&=\mu Y(X+1)+1,&H&=4A-5,&\Delta&=A^2-1,\\
 V_0&=X(\mu Y)^2,&E&=V_0(V_0+1)k^2+1,\\
 R&=\bigl[(q-1)(q-F)+\alpha\bigr](q^2-1)\\
  &\qquad +(M_C+qM_F)J,\\
 e&=k-R-1,&C&=q-F-\alpha,&U&=ic^2,\\
 D_f&=1+\Delta U^2E,&Q&=\Delta^2U^2E,&\delta&=Q-1,\\
 b&=c(T+\alpha)-R,&V&=D_fb-c,&D_y&=QV^2-1,\\
 K_+&=X(3c-2X)+2.
\end{align*}
The ordinary chart uses
\[
 B=(K_0+w)C-F,\quad h=q-1,\quad P=q^3(q-1),\quad D_{\rm low}=0,
\]
and the direct chart uses
\[
 B=(qK_0+X)C-qF,\quad h=q(q-1),\quad
 P=q^4(q-1),\quad D_{\rm low}=(q-1)X.
\]
Select a center
\[
 C_g=0,\qquad C_g=(A+1)UE,\qquad\text{or}\qquad C_g=AUE,
\]
respectively. Finally set
\begin{align*}
 M&=P\mu^2H,\\
 Z&=P\mu^2\{H(n+V+C_g-bB)+K_+\}\\
  &\qquad +H\{\mu^2(hY-q^3B-D_{\rm low})-Pe\},\\
 s&=\delta Z^2,\qquad u=\delta M^2ED_f,\qquad v=M^2D_y.
\end{align*}
The single equation, for any of the six choices, is
\[
 \boxed{N=(s-u-v)^2-4uv=0.}
 \tag{16}\label{eq:positiveN}
\]
All abbreviations use only addition, subtraction, multiplication, and fixed
nonnegative integer exponents. The literal JSON circuits expand every one
of them; no square root or auxiliary predicate occurs in $N$.

\subsection{The auxiliary coordinate is a sufficient transport bound}
The inherited coprime recovery, before any transport sign conclusion, gives
positive integers $\tau,f,y$ with
\[
 E=\tau^2,\quad D_f=f^2,\quad I=U\tau,\quad
 Q=(\Delta I)^2,\quad QV^2-(Q-1)y^2=1.
 \tag{17}\label{eq:recovered}
\]
For clarity, here $I$ is a proof-only recovered integer, not an additional
unknown. The coprime step is valid because $D_f\equiv1\pmod E$.
The two rational radicands $ED_f$ and $D_y/\delta$ cannot coincide:
after multiplying by $\delta$ and reducing modulo $E$, equality would give
$0\equiv-1$. Thus the rational signed-sum argument makes both roots
rational; square-class coprimality separates $E$ and $D_f$, and the
retained adjacent-coefficient descent makes $y$ integral. This order does
not presuppose the later projection or the sign of $B$.

Write $\psi_A(m)$ for the Pell ordinate at base $A$.
The retained uniform folding and first-index arguments give
$c=\psi_A(p)$ for an odd positive $p$, together with
$c>8q^5$, $c>2X$, $H>3X$, and $K_+>0$.
The native mask bound
$R<q^4+(\alpha+2)q^2<c\alpha/4$ then gives
\[
 b=c(T+\alpha)-R>c.
\]
Since $c\ge3$, the odd index $p$ is at least three, so
$c\ge4A^2-1=4\Delta+3$.

\begin{lemma}[Fifth-index gap after fusion]
Under the recovered conditions above,
\[
 b>8Q,\qquad g:=\tau f<Q,\qquad 0<y-V<b/4.
\]
\end{lemma}
\begin{proof}
From $b>c$,
\[
 V=D_fb-c>(D_f-1)c=\Delta I^2c>4Q.
\]
Put $S=\Delta I$, so $Q=S^2$. Pell classification applied to
$(SV)^2-(S^2-1)y^2=1$ gives $SV=\chi_S(\nu)$ with positive odd
$\nu$. The values at $\nu=1,3$ give $V=1,4Q-3$, both excluded.
Thus $\nu\ge5$ and
\[
 V\ge16Q^2-20Q+5>8Q^2.
\]
Since $D_f<Q$, this yields $b=(V+c)/D_f>8Q$.
Also $f<AI$, so $g<AUE<Q$.
Finally $Q-1>2D_f$ and
\[
 0<y-V=\frac{V^2-1}{(Q-1)(y+V)}
       <\frac{V}{2(Q-1)}<\frac b4.
\]
Each inequality precedes the transport sign recovery.
\end{proof}

For explanation only, divide the carrier by $M$ and put $W=1+hb$.
The exact identity is
\[
 \frac ZM=n+V+C_g+\frac{K_+}{H}+\frac{Y}{q^3}
          -\frac{WB}{h}-\frac e{\mu^2}-d,
 \tag{18}\label{eq:carrierpositive}
\]
where $d=0$ ordinarily and $d=X/q^4$ directly. At a norm zero,
$Z/M=\varepsilon_g g+\varepsilon_y y$.
Since $Q>X$, we have $g+d<2Q<b/4$. Also
$b-e=c(T+\alpha)-k+1>0$, so $e/\mu^2<b/4$ even when $e$ is negative.
If $B\le-1$, solving \eqref{eq:carrierpositive} for $n$ gives
\[
 n<\frac{3b}{4}+\frac{WB}{h}
   \le\frac{3b}{4}-b-\frac1h<0.
\]
This contradicts the domain of $n$. A positive center $C_g$ only decreases
the right side. Thus all six equations force $B\ge0$ using the existing
coordinate $b$ itself; the old separately computed transport majorant is
unnecessary. The cheapest pair can also discard the center.

\subsection{Completeness and the consistency theorem}
After $B\ge0$ is established, the unchanged source arguments recover
$C>0$, $3q<R<q^4-2$, the exact index $p=R$, denominator divisibilities,
and $X=2^R$. They reconstruct the same unscaled native source witness;
the scaled strong-Pell witness is not divided and reused without proof.
The retained native/controller/refutation calibration therefore applies.

Conversely, start with the parent's canonical data supplied by a finite
refutation. Choose a fresh strong Pell period with
$c^2\tau\mid I$, which is possible by the finite order of the
determinant-one Pell matrix modulo $c^2\tau$. Put
$i=I/(c^2\tau)>0$. The retained auxiliary-completeness construction gives
positive $T$, the auxiliary value $V$, and a gap $u_*=y-V>4I^2$.
It also gives
\[
 \kappa=K_+/H<c<I,\quad s_Y=Y/q^3<Y<c<I,\quad
 UE=I\tau<I^2.
\]
First work on the ordinary chart, with
$B_t=(K_0+w)C-F$ and $W=1+(q-1)b$.
The inherited canonical divisibilities make $\kappa$, $s_Y$,
$t_*=B_t/(q-1)\ge0$, and $r_*=e/\mu^2\ge0$ integers. Choose
\[
 n=g-C_g+u_*-\kappa-s_Y+Wt_*+r_*.
\]
For the uncentered choice $g>2c>\kappa+s_Y$ proves $n>0$.
For the centered choices, $0<C_g-g<2UE$ and the total negative loss
is less than $2I^2+2I<4I^2<u_*$. The norm then vanishes.
The direct chart takes $x=X-q>0$ and adds the positive integral correction
$(W_{\rm direct}-W_{\rm ordinary})t_*+X/q^4$ to $n$.
Here the inherited canonical power divisibility gives $q^4\mid X$.
Thus every finite refutation supplies positive zeros in both charts.

It follows that, for every polynomial $N$ in the table,
\[
 \ZFC\vdash
 [\forall\mathbf z\in\Z_{>0}^{10}\;N(\mathbf z)\ne0]
 \longleftrightarrow\Con(\ZFC).
\]
The complete new proof, with the precise inherited theorem interfaces,
is \file{positive/PROOF.md}.

\subsection{Exact degree, rather than a syntactic upper estimate}
Let $q_0=\beta J$, $k_0=\eta+\zeta$, and $x_0=q_0+x$.
Uniform formal scaling of all ten variables gives the following nonzero
highest homogeneous parts for the uncentered pair:
\begin{align*}
 N_{288,\mathrm{top}}
 &=-2^{10}q_0^{152}w^{30}Y^{62}i^{10}k_0^{32}(q_0-F)^2,\\
 N_{262,\mathrm{top}}
 &=-2^{10}q_0^{126}x_0^{30}Y^{62}i^{10}k_0^{32}(q_0-F)^2.
\end{align*}
They are nonzero for the actual positive fixed $\beta$.
In the formal Laurent-series field, the rational remainder in
$Z/M=V+C_g+\mathcal E$ has degree at most six. The centered radical
defects have degrees at least seventeen, so the changed transport cannot
alter their leading terms. Pairing the two signs of each radical yields
the exact centered degrees in the table. The proof provides the complete
leading forms and global degree bound; exact one-variable specialization
checks are supplementary transcription controls, not the sole degree proof.
